% !TEX output_directory = ../publica
\documentclass[twoside, 10pt, a4paper]{article}
\usepackage[utf8]{inputenc}
\usepackage[T1]{fontenc}
\usepackage[portuguese]{babel}

% --- FONTE PADRÃO (Estilo Didático Encorpado) ---
\usepackage{helvet} 
\renewcommand{\familydefault}{\sfdefault}

% --- GEOMETRIA E ESPAÇAMENTO ---
\usepackage[top=1.6cm, bottom=1.5cm, left=0.9cm, right=0.9cm, headheight=22pt, headsep=0.4cm]{geometry}
\usepackage{microtype} 
\usepackage{setspace}
\setstretch{1.0}
\usepackage{parskip}
\setlength{\parskip}{2pt}
\setlength{\parindent}{0pt}

\newlength{\espacoPosTitulo}\setlength{\espacoPosTitulo}{0.3cm}
\newlength{\espacoPreQuestao}\setlength{\espacoPreQuestao}{0.1cm}
\newlength{\espacoQuestao}\setlength{\espacoQuestao}{0.2cm}
\newlength{\espacoAlt}\setlength{\espacoAlt}{0.2cm}

\usepackage{enumitem}
\setlist{itemsep=0.4cm, topsep=0.1cm, parsep=0pt, partopsep=0pt}
\newcommand{\larguraTikz}{0.85\linewidth} 

% --- MATEMÁTICA E CORES ---
\usepackage{amsmath, amsfonts, amssymb, nccmath, mathastext}
\usepackage{xcolor}
\definecolor{indigo}{HTML}{4F46E5}
\definecolor{CorTitUm}{HTML}{FF6F61}
\definecolor{CorTitDois}{HTML}{FF6F61}
\definecolor{CorTitTres}{HTML}{658894}
\definecolor{CorLinha}{HTML}{B0B0B0} 
\definecolor{metal}{HTML}{7F8C8D}
\definecolor{vidro}{HTML}{3498DB}
\definecolor{ouro}{HTML}{F1C40F}
\definecolor{concreto}{HTML}{95A5A6}
\definecolor{madeira}{HTML}{D35400}
\definecolor{CinzaEscuro}{HTML}{4A4A4A} 
\definecolor{CinzaMedio}{HTML}{848484} 
\definecolor{Cinzablack}{HTML}{353535} 

% --- MINI-CABEÇALHO E RODAPÉ AUTORAL (TWOSIDE) ---
\usepackage{fancyhdr}
\pagestyle{fancy}
\fancyhf{} 
\renewcommand{\headrulewidth}{0.3pt} 
\renewcommand{\footrulewidth}{0.3pt}
\fancyhead[LE]{\small\sffamily\bfseries\color{gray!75} MÁQUINAS SIMPLES}
\fancyhead[RO]{\small\sffamily\color{gray!75} \texttt{www.profraphaelpaulino.com.br}}
\fancyfoot[LE]{\small \textbf{\thepage}}
\fancyfoot[RO]{\small \textbf{\thepage}}

% --- CAIXAS DE DESTAQUE ---
\usepackage[most]{tcolorbox}
\newtcolorbox{boxresponda}{colback=white, colframe=gray!45, boxrule=0pt, leftrule=1.7pt, sharp corners, enlarge left by=0.4cm, width=\linewidth-0.4cm, left=8pt, right=0pt, top=2pt, bottom=2pt, fontupper=\small}
\definecolor{CorDicaBorda}{HTML}{17c1be} 
\definecolor{CorDicaFundo}{HTML}{f3fbfb} 
\newtcolorbox{boxdica}{colback=CorDicaFundo, colframe=CorDicaBorda, boxrule=0pt, leftrule=2.5pt, sharp corners, width=\linewidth, left=8pt, right=8pt, top=6pt, bottom=6pt, halign=center, fontupper=\small}

% --- TÍTULOS ---
\usepackage{titlesec}
\titleformat{\section}{\color{black}\large\bfseries}{\thesection}{0em}{\vspace{0.1cm}\titlerule[0.8pt]}
\titlespacing*{\section}{0pt}{10pt}{6pt} 
\titleformat{\subsubsection}[runin]{\color{black}\bfseries}{}{0em}{}
\titlespacing*{\subsubsection}{0pt}{0pt}{0.15cm}

% --- COLUNAS E GRÁFICOS ---
\usepackage{multicol}
\setlength{\columnsep}{1cm}
\setlength{\columnseprule}{0.5pt}
\def\columnseprulecolor{\color{CorLinha}}
\usepackage{adjustbox, tikz, pgfplots, circuitikz}
\usetikzlibrary{shadows, shapes.misc, positioning, calc}
\pgfplotsset{compat=1.18}

\begin{document}
\thispagestyle{empty}
\color{Cinzablack}
\vspace*{-1.8cm}

% --- CABEÇALHO PRINCIPAL AUTORAL (Apenas 1ª página) ---
\noindent
\begin{minipage}{\textwidth}
    \fontfamily{phv}\selectfont
    \noindent
    \begin{minipage}[t]{0.70\linewidth}
        {\Large \bfseries \MakeUppercase{Caderno de Atividades Suplementares}}\\[0.1cm]
        {\large \textmd{\textcolor{CinzaEscuro}{Unidade: MÁQUINAS SIMPLES}}}
    \end{minipage}%
    \begin{minipage}[t]{0.30\linewidth}
        \raggedleft
        \small \textbf{Prof. Raphael Paulino}\\
        \textsf{\small \textcolor{indigo}{\texttt{profraphaelpaulino.com.br}}}
    \end{minipage}
    
    \vspace{0.15cm}
    \noindent\rule{\linewidth}{1.2pt}
    \vspace{-0.1cm}
    \footnotesize\textsf{\textbf{ÁREA:} FÍSICA \hfill \textbf{NÍVEL:} EF II (7º ANO)}
    \vspace{0.1cm}
    \noindent\rule{\linewidth}{0.4pt}
\end{minipage}

\par\vspace{0.3cm} 
\raggedcolumns
\begin{multicols*}{2}
\vspace{0.1cm}

\subsection*{Fundamentos e Evolução Histórica}
\vspace{-0.2cm}\noindent\rule{\linewidth}{1.5pt} 
\par\vspace{\espacoPosTitulo}
\subsubsection*{1.} A história da humanidade está diretamente ligada à invenção de dispositivos que facilitam o trabalho. Desde a construção das pirâmides do Egito até os aquedutos romanos, a manipulação de forças foi essencial. Sobre o conceito físico de trabalho e a atuação das máquinas simples, analise as afirmativas a seguir e classifique-as como Verdadeiras (\textbf{V}) ou Falsas (\textbf{F}):
\begin{enumerate}[label=\Roman*., itemsep=\espacoAlt]
    \item As máquinas simples reduzem a quantidade total de trabalho físico realizado para deslocar um objeto.
    \item Uma máquina simples pode alterar a direção da força aplicada, facilitando a execução de uma tarefa.
    \item O uso de um plano inclinado diminui a força necessária para elevar uma carga, mas exige que ela percorra uma distância maior.
\end{enumerate}
\par\vspace{\espacoQuestao}

\noindent\textcolor{CorLinha}{\rule{\linewidth}{0.5pt}} 
\par\vspace{\espacoPreQuestao}
\subsubsection*{2.} Durante a Idade Média, o desenvolvimento de fortificações e armamentos de cerco, como as catapultas, exigiu um refinamento empírico da mecânica. Do ponto de vista da Física, explique qual é a principal vantagem mecânica obtida ao se utilizar um dispositivo que atua como multiplicador de força.
\par\vspace{\espacoQuestao}

\noindent\textcolor{CorLinha}{\rule{\linewidth}{0.5pt}} 
\par\vspace{\espacoPreQuestao}
\subsubsection*{3.} O esquema a seguir ilustra o princípio de funcionamento de uma catapulta clássica, onde uma força é aplicada em uma extremidade para lançar um projétil na outra.
\begin{center}
% ==========================================
% CONTROLE INDIVIDUAL DESTA IMAGEM:
% 1. CORTAR BORDAS: Mude os zeros do trim (Esquerda, Baixo, Direita, Topo). Ex: trim=1cm 0cm 1cm 0cm
% 2. TAMANHO: Para encolher só esta imagem, troque \larguraTikz por um valor (Ex: 0.5\linewidth)
% ==========================================
\begin{adjustbox}{trim=0cm 0cm 0cm 0cm, clip, max width=\larguraTikz}
\begin{tikzpicture}
% --- APLICANDO DROP SHADOW, CORES HEX E BORDAS ARREDONDADAS ---
\definecolor{madeira_clara}{HTML}{8B5A2B}
\definecolor{pedra}{HTML}{95A5A6}

% LAYER 1: Fundo (Chão)
\draw[drop shadow={opacity=0.15}, fill=gray!15, rounded corners=2pt, draw=none] (-3.5,-0.5) rectangle (3.5,0);

% LAYER 2: Ponto de apoio
\draw[drop shadow={opacity=0.2}, fill=metal, rounded corners=1pt] (-0.5,0) -- (0.5,0) -- (0,1) -- cycle;

% LAYER 3: Haste da alavanca
\draw[drop shadow={opacity=0.2}, fill=madeira_clara, rounded corners=2pt, rotate around={15:(0,1)}] (-2.8,0.9) rectangle (2.8,1.1);

% LAYER 4: Carga (projétil)
\draw[drop shadow={opacity=0.2}, fill=pedra, rounded corners=3pt, rotate around={15:(0,1)}] (-2.6, 1.1) rectangle (-1.6, 1.9);
\node[rotate=15] at (-2.3, 2.1) [fill=white, inner sep=1pt] {\textbf{Projétil}};

% LAYER 5: Força aplicada
\draw[->, line width=1.5pt, draw=black] (2.2, 2.2) -- (2.2, 1) node[midway, right, fill=white, inner sep=2pt] {\textbf{F}};
\node at (0, -0.25) [fill=gray!15, inner sep=2pt] {\textbf{Apoio}};
\end{tikzpicture}
\end{adjustbox}
\end{center}
Considerando as posições da força motriz (\textbf{F}), do projétil (resistência) e do ponto de apoio geométrico, classifique o tipo de alavanca operada pelo mecanismo e justifique sua resposta.
\par\vspace{\espacoQuestao}

\noindent\textcolor{CorLinha}{\rule{\linewidth}{0.5pt}}
\par\vspace{\espacoPreQuestao}
\subsubsection*{4.} Um estudante do 7º ano, ao analisar as engrenagens da Revolução Industrial, escreveu em seu relatório: "Como as polias e alavancas multiplicam a força, podemos criar uma máquina simples que aumente simultaneamente a força aplicada e a velocidade de deslocamento de uma carga pesada, sem precisar de energia externa adicional."
\begin{boxresponda}
\textbf{Responda:} Diagnostique o erro físico na afirmação do estudante. Explique a relação de compensação obrigatória entre força e distância quando utilizamos uma máquina simples.
\end{boxresponda}
\par\vspace{\espacoQuestao}

\noindent\textcolor{CorLinha}{\rule{\linewidth}{0.5pt}} 
\par\vspace{\espacoPreQuestao}
\subsubsection*{5.} (Estilo UNESP - Adaptada) As ferramentas utilizadas diariamente em residências e oficinas são, na sua grande maioria, aplicações diretas dos princípios físicos das máquinas simples estabelecidos desde a Antiguidade. Considere os seguintes instrumentos:
1. Uma tesoura de costura.
2. Um abridor de garrafas.
3. Uma maçaneta de porta convencional.
Os princípios físicos predominantes em cada um desses instrumentos correspondem, respectivamente, a:
\begin{enumerate}[label=\alph*), itemsep=\espacoAlt]
    \item Alavanca interfixa, plano inclinado, roda e eixo.
    \item Alavanca interpotente, alavanca interfixa, polia.
    \item Alavanca interfixa, alavanca inter-resistente, roda e eixo.
    \item Plano inclinado, alavanca inter-resistente, polia fixa.
    \item Alavanca interpotente, plano inclinado, roda e eixo.
\end{enumerate}
\par\vspace{\espacoQuestao}

\subsection*{Princípios Físicos (Alavancas, Polias e Planos)}
\vspace{-0.2cm}\noindent\rule{\linewidth}{1.5pt} 
\par\vspace{\espacoPosTitulo}
\subsubsection*{6.} Para manter o sistema em repouso perfeito na horizontal, um professor de física posicionou dois blocos maciços sobre uma prancha articulada, conforme o arranjo geométrico ilustrado a seguir.
\begin{center}
% ==========================================
% CONTROLE INDIVIDUAL DESTA IMAGEM:
% ==========================================
\begin{adjustbox}{trim=0cm 0cm 0cm 0cm, clip, max width=\larguraTikz}
\begin{tikzpicture}
% --- APLICANDO DROP SHADOW, CORES HEX E BORDAS ARREDONDADAS ---
\definecolor{base}{HTML}{2C3E50}
\definecolor{prancha}{HTML}{F39C12}
\definecolor{blocoA}{HTML}{2980B9}
\definecolor{blocoB}{HTML}{C0392B}

% LAYER 1: Chão e linhas de cota
\draw[fill=gray!10, draw=none] (-4,-1) rectangle (4, -0.5);
\draw[thick, color=gray!50] (-4,-0.5) -- (4,-0.5);

% LAYER 2: Apoio central
\draw[drop shadow={opacity=0.2}, fill=base, rounded corners=1pt] (-0.5,-0.5) -- (0.5,-0.5) -- (0,0.5) -- cycle;

% LAYER 3: Prancha
\draw[drop shadow={opacity=0.2}, fill=prancha, rounded corners=2pt] (-3.5,0.5) rectangle (3.5,0.7);

% LAYER 4: Blocos com massas
\draw[drop shadow={opacity=0.2}, fill=blocoA, rounded corners=2pt] (-3,0.7) rectangle (-2,1.7) node[pos=.5, text=white] {\textbf{60 kg}};
\draw[drop shadow={opacity=0.2}, fill=blocoB, rounded corners=2pt] (1.5,0.7) rectangle (2.5,1.7) node[pos=.5, text=white] {\textbf{m}};

% LAYER 5: Cotas dimensionais com máscara de fundo branco
\draw[<->, >=stealth, thick] (-2.5, 2.3) -- (0, 2.3) node[midway, fill=white, inner sep=2pt] {\textbf{2,0 m}};
\draw[<->, >=stealth, thick] (0, 2.3) -- (2, 2.3) node[midway, fill=white, inner sep=2pt] {\textbf{3,0 m}};
\draw[dashed, color=gray!80] (-2.5, 1.7) -- (-2.5, 2.5);
\draw[dashed, color=gray!80] (0, 0.5) -- (0, 2.5);
\draw[dashed, color=gray!80] (2, 1.7) -- (2, 2.5);
\end{tikzpicture}
\end{adjustbox}
\end{center}
Determine o valor numérico da massa \textbf{m} do bloco vermelho para que a condição de equilíbrio estático rotacional seja satisfeita.
\par\vspace{\espacoQuestao}

\noindent\textcolor{CorLinha}{\rule{\linewidth}{0.5pt}} 
\par\vspace{\espacoPreQuestao}
\subsubsection*{7.} Em um canteiro de obras, um operário precisa erguer uma caixa de ferramentas utilizando um sistema de içamento acoplado ao teto do andaime, como esquematizado na figura abaixo.
\begin{center}
\begin{adjustbox}{trim=0cm 0cm 0cm 0cm, clip, max width=\larguraTikz}
\begin{tikzpicture}
% --- APLICANDO DROP SHADOW, CORES HEX E BORDAS ARREDONDADAS ---
\definecolor{teto}{HTML}{7F8C8D}
\definecolor{corda}{HTML}{34495E}
\definecolor{polia}{HTML}{BDC3C7}
\definecolor{carga}{HTML}{8E44AD}

% LAYER 1: Teto
\draw[fill=teto, drop shadow={opacity=0.2}] (-2,4) rectangle (2,4.3);

% LAYER 2: Polias e suportes
\draw[fill=polia, drop shadow={opacity=0.15}] (1,3) circle (0.4);
\draw[fill=gray!80] (1,3) circle (0.1);
\draw[thick, color=teto] (1,3) -- (1,4); 

\draw[fill=polia, drop shadow={opacity=0.15}] (-1,2) circle (0.4);
\draw[fill=gray!80] (-1,2) circle (0.1);

% LAYER 3: Cabeamento (Cordas)
\draw[thick, color=corda] (-1.4,4) -- (-1.4,2); 
\draw[thick, color=corda] (-0.6,2) -- (0.6,3); 
\draw[thick, color=corda, ->, >=stealth] (1.4,3) -- (1.4,1) node[below, fill=white, inner sep=2pt] {\textbf{F}};

% LAYER 4: Carga
\draw[thick, color=teto] (-1,2) -- (-1,1.2); 
\draw[fill=carga, drop shadow={opacity=0.2}, rounded corners=2pt] (-1.6,0.2) rectangle (-0.4,1.2) node[pos=.5, text=white] {\textbf{800 N}};
\end{tikzpicture}
\end{adjustbox}
\end{center}
Com base na associação de roldanas empregada, calcule o valor da força motriz ideal \textbf{F} que o operário deve aplicar verticalmente para baixo, a fim de içar a carga com velocidade constante.
\par\vspace{\espacoQuestao}

\noindent\textcolor{CorLinha}{\rule{\linewidth}{0.5pt}} 
\par\vspace{\espacoPreQuestao}
\subsubsection*{8.} A alavanca é composta basicamente por uma barra rígida e um ponto de apoio. Dependendo da posição do ponto de apoio em relação à força potente e à força resistente, ela recebe uma classificação específica. Identifique e descreva a classificação física (interfixa, interpotente ou inter-resistente) de cada um dos instrumentos listados:
\begin{enumerate}[label=\Roman*., itemsep=\espacoAlt]
    \item Um quebra-nozes mecânico.
    \item Uma pinça de sobrancelha.
    \item Uma gangorra de parque.
\end{enumerate}
\par\vspace{\espacoQuestao}

\noindent\textcolor{CorLinha}{\rule{\linewidth}{0.5pt}} 
\par\vspace{\espacoPreQuestao}
\subsubsection*{9.} A rampa de acesso ao setor de logística de um armazém foi projetada com medidas específicas para minimizar o esforço dos funcionários, conforme a representação geométrica a seguir.
\begin{center}
\begin{adjustbox}{trim=0cm 0cm 0cm 0cm, clip, max width=\larguraTikz}
\begin{tikzpicture}
% --- APLICANDO DROP SHADOW E BORDAS ARREDONDADAS ---
\definecolor{rampa}{HTML}{BDC3C7}
\definecolor{caixa}{HTML}{D35400}

% LAYER 1: Estrutura da Rampa
\draw[fill=rampa, drop shadow={opacity=0.15}, draw=gray!70] (0,0) -- (5,0) -- (5,2.5) -- cycle;
\draw[thick, color=concreto] (-1,0) -- (6,0);

% LAYER 2: Cotas dimensionais
\draw[thick, color=black] (5,0) -- (5,2.5) node[midway, right, fill=white, inner sep=2pt] {\textbf{2,0 m}};
\draw[thick, color=black] (0,0) -- (5,2.5) node[pos=0.25, above left, fill=white, inner sep=2pt, sloped] {\textbf{10,0 m}};

% LAYER 3: Objeto na rampa
\draw[fill=caixa, drop shadow={opacity=0.2}, rounded corners=1pt, rotate around={26.5:(2,1)}] (2,1) rectangle (2.8,1.6) node[pos=.5, text=white, rotate=26.5] {\textbf{500 N}};

% LAYER 4: Vetor Força Motriz
\draw[->, >=stealth, line width=1.2pt, rotate around={26.5:(2,1)}] (2.8, 1.3) -- (4, 1.3) node[above, sloped, fill=white, inner sep=1pt] {\textbf{F}};
\end{tikzpicture}
\end{adjustbox}
\end{center}
Sem considerar as perdas por atrito, aplique o princípio da vantagem mecânica do plano inclinado e determine o valor da força \textbf{F} necessária para empurrar o objeto uniformemente até o topo.
\par\vspace{\espacoQuestao}

\noindent\textcolor{CorLinha}{\rule{\linewidth}{0.5pt}} 
\par\vspace{\espacoPreQuestao}
\subsubsection*{10.} (Estilo ENEM - Adaptada) A combinação de uma roda rigidamente presa a um eixo concêntrico constitui uma máquina simples cujo objetivo é transferir torque e amplificar forças rotacionais. Nas aplicações práticas modernas, essa mecânica é vital. Sobre a "roda e eixo", julgue as afirmativas a seguir assinalando Verdadeiro (\textbf{V}) ou Falso (\textbf{F}):
\begin{enumerate}[label=\Roman*., itemsep=\espacoAlt]
    \item Quando aplicamos força na extremidade do volante de um automóvel (roda), a força transferida para a barra de direção (eixo) é menor do que a aplicada.
    \item O diâmetro maior da maçaneta de uma porta em relação ao seu eixo interno garante uma vantagem mecânica que facilita o destravamento da fechadura.
    \item Se a força motriz for aplicada no eixo em vez de na roda (como nos pneus traseiros de um carro), haverá um ganho de velocidade periférica em sacrifício da força.
\end{enumerate}
\par\vspace{\espacoQuestao}

\noindent\textcolor{CorLinha}{\rule{\linewidth}{0.5pt}} 
\par\vspace{\espacoPreQuestao}
\subsubsection*{11.} Sistemas compostos por múltiplas polias móveis, conhecidos como talhas exponenciais, são projetados para levantar cargas massivas em estaleiros e portos. O sistema esquematizado opera içando um contêiner sob tensão contínua.
\begin{center}
\begin{adjustbox}{trim=0cm 0cm 0cm 0cm, clip, max width=\larguraTikz}
\begin{tikzpicture}
% --- APLICANDO DROP SHADOW, CORES HEX E BORDAS ARREDONDADAS ---
\definecolor{tetopolia}{HTML}{34495E}
\definecolor{poliaF}{HTML}{95A5A6}
\definecolor{poliaM}{HTML}{F1C40F}
\definecolor{bloco}{HTML}{E74C3C}

% LAYER 1: Teto suporte
\draw[fill=tetopolia, drop shadow={opacity=0.2}] (-2,5) rectangle (3,5.3);

% LAYER 2: Polia Fixa
\draw[fill=poliaF, drop shadow={opacity=0.15}] (2,4) circle (0.4);
\draw[fill=white] (2,4) circle (0.1);
\draw[thick] (2,4) -- (2,5);

% LAYER 3: Polias Móveis
\draw[fill=poliaM, drop shadow={opacity=0.15}] (0.5,3) circle (0.4);
\draw[fill=white] (0.5,3) circle (0.1);

\draw[fill=poliaM, drop shadow={opacity=0.15}] (-1,2) circle (0.4);
\draw[fill=white] (-1,2) circle (0.1);

% LAYER 4: Cordas (Talha Exponencial)
\draw[thick] (-1.4,5) -- (-1.4,2); 
\draw[thick] (-0.6,2) -- (-0.6,3); 
\draw[thick] (-0.6,3) -- (0.5,3); 
\draw[thick] (0.1,5) -- (0.1,3);
\draw[thick] (0.9,3) -- (0.9,4);
\draw[thick] (0.9,4) -- (2,4);
\draw[thick] (1.6,5) -- (1.6,4);
\draw[thick, ->, >=stealth] (2.4,4) -- (2.4,2) node[below, fill=white, inner sep=2pt] {\textbf{F}};

% LAYER 5: Carga pesada
\draw[thick] (-1,2) -- (-1,1);
\draw[fill=bloco, drop shadow={opacity=0.2}, rounded corners=2pt] (-1.8,0) rectangle (-0.2,1) node[pos=.5, text=white] {\textbf{1600 N}};
\end{tikzpicture}
\end{adjustbox}
\end{center}
Neste arranjo contendo exatamente duas polias móveis e uma fixa, calcule a intensidade da força \textbf{F} necessária para equilibrar a carga.
\par\vspace{\espacoQuestao}

\noindent\textcolor{CorLinha}{\rule{\linewidth}{0.5pt}}
\par\vspace{\espacoPreQuestao}
\subsubsection*{12.} Uma das aplicações mais eficientes da associação entre rodas e eixos pode ser observada no sistema de marchas (coroa e catraca) de uma bicicleta moderna. Um aluno, ao tentar explicar o funcionamento, disse que "ao passar para uma catraca traseira maior, o ciclista pedala mais rápido e com mais força na roda".
\begin{boxresponda}
\textbf{Responda:} Refute essa afirmação provando fisicamente o que realmente ocorre com a velocidade da bicicleta e com a força exigida do ciclista quando a corrente é mudada para uma catraca de raio muito maior.
\end{boxresponda}
\par\vspace{\espacoQuestao}
\noindent\textcolor{CorLinha}{\rule{\linewidth}{0.5pt}} 
\par\vspace{\espacoPreQuestao}
\subsubsection*{13.} A necessidade de elevar materiais pesados em construções civis popularizou o uso de polias. Um mestre de obras montou o esquema ilustrado abaixo, associando uma roldana fixa e uma móvel para erguer um bloco de cimento.
\begin{center}
% ==========================================
% CONTROLE INDIVIDUAL DESTA IMAGEM:
% ==========================================
\begin{adjustbox}{trim=0cm 0cm 0cm 0cm, clip, max width=\larguraTikz}
\begin{tikzpicture}
% --- APLICANDO DROP SHADOW, CORES HEX E BORDAS ARREDONDADAS ---
\definecolor{viga}{HTML}{2C3E50}
\definecolor{roldana}{HTML}{7F8C8D}
\definecolor{cimento}{HTML}{95A5A6}
\definecolor{cordame}{HTML}{E67E22}

% LAYER 1: Viga superior
\draw[fill=viga, drop shadow={opacity=0.2}] (-2,4) rectangle (2,4.3);

% LAYER 2: Roldanas (Fixa e Móvel)
\draw[fill=roldana, drop shadow={opacity=0.15}] (1,3) circle (0.4);
\draw[fill=white] (1,3) circle (0.1);
\draw[thick, color=viga] (1,3) -- (1,4); % Eixo da fixa

\draw[fill=roldana, drop shadow={opacity=0.15}] (-1,2) circle (0.4);
\draw[fill=white] (-1,2) circle (0.1);

% LAYER 3: Cordame
\draw[thick, color=cordame] (-1.4,4) -- (-1.4,2); 
\draw[thick, color=cordame] (-0.6,2) -- (0.6,3); 
\draw[thick, color=cordame, ->, >=stealth] (1.4,3) -- (1.4,1) node[below, fill=white, inner sep=2pt] {\textbf{F}};

% LAYER 4: Bloco de cimento
\draw[thick, color=viga] (-1,2) -- (-1,1.2);
\draw[fill=cimento, drop shadow={opacity=0.2}, rounded corners=2pt] (-1.5,0.2) rectangle (-0.5,1.2) node[pos=.5, text=black] {\textbf{500 N}};
\end{tikzpicture}
\end{adjustbox}
\end{center}
Determine o valor da força motriz \textbf{F} que o operário deverá aplicar na extremidade livre da corda para manter o bloco de cimento em equilíbrio estático no ar.
\par\vspace{\espacoQuestao}

\noindent\textcolor{CorLinha}{\rule{\linewidth}{0.5pt}} 
\par\vspace{\espacoPreQuestao}
\subsubsection*{14.} O uso de planos inclinados em estradas de montanha é uma aplicação direta da física para permitir que veículos alcancem grandes altitudes. Sobre a vantagem mecânica do plano inclinado, julgue as afirmativas a seguir com Verdadeiro (\textbf{V}) ou Falso (\textbf{F}):
\begin{enumerate}[label=\Roman*., itemsep=\espacoAlt]
    \item Quanto mais extensa (comprida) for a rampa para atingir uma mesma altura, menor será a força necessária para subir o veículo.
    \item O plano inclinado reduz o trabalho total realizado pelo motor do carro, pois encurta a distância do trajeto.
    \item A força motriz necessária para empurrar um objeto rampa acima é inversamente proporcional à inclinação; rampas mais suaves exigem menos força.
\end{enumerate}
\par\vspace{\espacoQuestao}

\noindent\textcolor{CorLinha}{\rule{\linewidth}{0.5pt}} 
\par\vspace{\espacoPreQuestao}
\subsubsection*{15.} O carrinho de mão é uma das invenções mais geniais da mecânica clássica, operando como uma alavanca que facilita o transporte de cargas pesadas por uma única pessoa. O diagrama abaixo detalha as posições do eixo da roda, do centro de massa da carga e dos puxadores.
\begin{center}
\begin{adjustbox}{trim=0cm 0cm 0cm 0cm, clip, max width=\larguraTikz}
\begin{tikzpicture}
% --- APLICANDO DROP SHADOW, CORES HEX E BORDAS ARREDONDADAS ---
\definecolor{metal_carrinho}{HTML}{34495E}
\definecolor{roda}{HTML}{2C3E50}
\definecolor{cargacaixa}{HTML}{D35400}

% LAYER 1: Solo e Roda (Ponto de Apoio)
\draw[thick, color=gray!50] (-1,0) -- (6,0);
\draw[fill=roda, drop shadow={opacity=0.2}] (0,0.5) circle (0.5);
\draw[fill=white] (0,0.5) circle (0.1) node[below=4pt, fill=white, inner sep=1pt] {\textbf{A}};

% LAYER 2: Estrutura do carrinho
\draw[thick, line width=2pt, color=metal_carrinho] (0,0.5) -- (4,1.5);
\draw[thick, line width=2pt, color=metal_carrinho] (4,1.5) -- (5,1.5) node[right] {\textbf{C}};

% LAYER 3: Carga
\draw[fill=cargacaixa, drop shadow={opacity=0.2}, rounded corners=3pt] (0.8, 1) -- (2.8, 1.5) -- (3, 2.5) -- (1, 2) -- cycle;
\node[fill=white, inner sep=1pt, rounded corners=1pt] at (2, 1.7) {\textbf{B (Carga)}};

% LAYER 4: Forças e Cotas
\draw[->, >=stealth, line width=1.2pt, color=black] (5,1.5) -- (5,2.5) node[above, fill=white, inner sep=1pt] {\textbf{F}};
\draw[->, >=stealth, line width=1.2pt, color=black] (2, 1.7) -- (2, 0.5) node[below, fill=white, inner sep=1pt] {\textbf{800 N}};

\draw[<->, >=stealth, dashed] (0, -0.5) -- (2, -0.5) node[midway, fill=white, inner sep=2pt] {\textbf{0,5 m}};
\draw[<->, >=stealth, dashed] (2, -0.5) -- (5, -0.5) node[midway, fill=white, inner sep=2pt] {\textbf{1,0 m}};
\draw[dashed, color=gray!80] (0,0) -- (0,-0.7);
\draw[dashed, color=gray!80] (2,0.5) -- (2,-0.7);
\draw[dashed, color=gray!80] (5,1.5) -- (5,-0.7);
\end{tikzpicture}
\end{adjustbox}
\end{center}
Sabendo que o carrinho atua como uma alavanca inter-resistente, calcule a força motriz \textbf{F} que o operador deve aplicar verticalmente para cima no ponto \textbf{C} para conseguir suspender o sistema e iniciar o transporte.
\par\vspace{\espacoQuestao}

\noindent\textcolor{CorLinha}{\rule{\linewidth}{0.5pt}} 
\par\vspace{\espacoPreQuestao}
\subsubsection*{16.} (Estilo FUVEST - Adaptada) Para carregar um cofre pesado de \textbf{1200 N} na carroceria de um caminhão que está a \textbf{1,5 m} de altura do solo, os entregadores utilizam uma prancha de madeira lisa atuando como plano inclinado. A prancha possui \textbf{4,5 m} de comprimento total. Desprezando qualquer força de atrito entre o cofre e a prancha, determine qual será a força constante, paralela à rampa, que os entregadores precisarão aplicar para empurrar o cofre de forma equilibrada até o topo.
\par\vspace{\espacoQuestao}

\noindent\textcolor{CorLinha}{\rule{\linewidth}{0.5pt}} 
\par\vspace{\espacoPreQuestao}
\subsubsection*{17.} Durante uma aula prática no parquinho da escola, um professor de física pediu que dois alunos de massas diferentes se sentassem em uma gangorra (alavanca interfixa) de forma a mantê-la perfeitamente na horizontal.
\begin{center}
\begin{adjustbox}{trim=0cm 0cm 0cm 0cm, clip, max width=\larguraTikz}
\begin{tikzpicture}
% --- APLICANDO DROP SHADOW, CORES HEX E BORDAS ARREDONDADAS ---
\definecolor{pedestal}{HTML}{16A085}
\definecolor{pranchamadeira}{HTML}{F39C12}
\definecolor{alunoA}{HTML}{8E44AD}
\definecolor{alunoB}{HTML}{2980B9}

% LAYER 1: Chão e Pedestal
\draw[fill=gray!15, draw=none] (-4.5,-0.5) rectangle (4.5,0);
\draw[thick, color=gray!50] (-4.5,0) -- (4.5,0);
\draw[fill=pedestal, drop shadow={opacity=0.2}, rounded corners=1pt] (-0.5,0) -- (0.5,0) -- (0,1) -- cycle;

% LAYER 2: Prancha da gangorra
\draw[fill=pranchamadeira, drop shadow={opacity=0.2}, rounded corners=2pt] (-4,1) rectangle (4,1.2);

% LAYER 3: Alunos (Blocos representativos)
\draw[fill=alunoA, drop shadow={opacity=0.2}, rounded corners=2pt] (-3.5,1.2) rectangle (-2.5,2.2) node[pos=.5, text=white] {\textbf{30 kg}};
\draw[fill=alunoB, drop shadow={opacity=0.2}, rounded corners=2pt] (1.5,1.2) rectangle (2.5,2.7) node[pos=.5, text=white] {\textbf{45 kg}};

% LAYER 4: Cotas de distância
\draw[<->, >=stealth, thick] (-3, 2.8) -- (0, 2.8) node[midway, fill=white, inner sep=2pt] {\textbf{1,8 m}};
\draw[<->, >=stealth, thick] (0, 2.8) -- (2, 2.8) node[midway, fill=white, inner sep=2pt] {\textbf{X}};
\draw[dashed, color=gray!80] (-3, 2.2) -- (-3, 3.0);
\draw[dashed, color=gray!80] (0, 1.2) -- (0, 3.0);
\draw[dashed, color=gray!80] (2, 2.7) -- (2, 3.0);
\end{tikzpicture}
\end{adjustbox}
\end{center}
Aplicando o princípio do equilíbrio das alavancas, determine a distância \textbf{X} que o aluno mais pesado (45 kg) deve ficar do eixo central para garantir que a gangorra não tombe.
\par\vspace{\espacoQuestao}

\noindent\textcolor{CorLinha}{\rule{\linewidth}{0.5pt}} 
\par\vspace{\espacoPreQuestao}
\subsubsection*{18.} Poços artesianos rústicos frequentemente utilizam um sistema de manivela acoplado a um tambor cilíndrico para içar baldes de água do fundo. Esse mecanismo é a clássica representação física da máquina simples conhecida como "roda e eixo". Se o tambor cilíndrico (onde a corda se enrola) possui um raio de \textbf{10 cm} e a haste da manivela possui um braço de alavanca de \textbf{40 cm}, explique fisicamente por que é muito mais fácil puxar um balde pesado de água girando a manivela do que puxando a corda diretamente com as mãos.
\par\vspace{\espacoQuestao}

\noindent\textcolor{CorLinha}{\rule{\linewidth}{0.5pt}}
\par\vspace{\espacoPreQuestao}
\subsubsection*{19.} (Estilo UNESP - Adaptada) Para remover o motor de um automóvel, um mecânico utiliza uma talha manual composta por uma associação de fios ideais contendo 1 polia fixa presa ao teto e 3 polias móveis encadeadas. Sabendo que o peso total do motor é de \textbf{1600 N}, as alternativas abaixo apresentam possíveis valores para a força ideal contínua que o mecânico precisa aplicar:
\begin{enumerate}[label=\alph*), itemsep=\espacoAlt]
    \item $ 1600 \text{ N} $
    \item $ 800 \text{ N} $
    \item $ 400 \text{ N} $
    \item $ 200 \text{ N} $
    \item $ 100 \text{ N} $
\end{enumerate}
\begin{boxresponda}
\textbf{Responda:} Identifique a alternativa correta e apresente o cálculo que refuta matematicamente as demais opções, provando por que a vantagem mecânica não é atingida nos valores errados.
\end{boxresponda}
\par\vspace{\espacoQuestao}

\subsection*{Soluções Práticas e Aplicações Cotidianas}
\vspace{-0.2cm}\noindent\rule{\linewidth}{1.5pt} 
\par\vspace{\espacoPosTitulo}
\subsubsection*{20.} A chave de fenda é uma ferramenta essencial que pode atuar como duas máquinas simples distintas dependendo de como é usada. Quando a utilizamos para forçar a abertura da tampa de uma lata de tinta, ela age como uma alavanca. Porém, quando a giramos para apertar um parafuso, ela atua pelo princípio de roda e eixo. Considere duas chaves de fenda de mesmo comprimento, mas uma possui o cabo (empunhadura) muito mais grosso (maior diâmetro) que a outra. Qual das duas chaves exigirá menos esforço físico do trabalhador para girar um parafuso enferrujado? Justifique sua escolha com base nos princípios físicos da vantagem mecânica.
\par\vspace{\espacoQuestao}

\noindent\textcolor{CorLinha}{\rule{\linewidth}{0.5pt}} 
\par\vspace{\espacoPreQuestao}
\subsubsection*{21.} O próprio corpo humano é um complexo sistema de máquinas simples. Ao segurarmos um objeto pesado com a mão e dobrarmos o cotovelo, o sistema osso-músculo atua exatamente como uma alavanca. O diagrama ilustra as forças em ação no antebraço.
\begin{center}
\begin{adjustbox}{trim=0cm 0cm 0cm 0cm, clip, max width=\larguraTikz}
\begin{tikzpicture}
% --- APLICANDO DROP SHADOW, CORES HEX E BORDAS ARREDONDADAS ---
\definecolor{pele}{HTML}{F5CBA7}
\definecolor{musculo}{HTML}{E74C3C}
\definecolor{pesomao}{HTML}{7F8C8D}

% LAYER 1: Braço superior e antebraço (Formas simplificadas)
\draw[fill=pele, drop shadow={opacity=0.15}, rounded corners=4pt] (-1,0) rectangle (0,4); % Braço
\draw[fill=pele, drop shadow={opacity=0.15}, rounded corners=4pt] (-0.5,-0.5) rectangle (5,0.5); % Antebraço

% LAYER 2: Articulação e Músculo
\draw[fill=white, thick] (0,0) circle (0.2) node[below=6pt, fill=white, inner sep=1pt] {\textbf{Articulação (Apoio)}};
\draw[fill=musculo, rounded corners=3pt, opacity=0.8] (-0.2, 3) to[out=-20, in=120] (1, 0.4) to[out=180, in=-60] (-0.8, 2) -- cycle;

% LAYER 3: Peso na mão
\draw[fill=pesomao, drop shadow={opacity=0.2}, rounded corners=2pt] (4.5,-0.2) rectangle (5.5,0.8) node[pos=.5, text=white] {\textbf{P}};

% LAYER 4: Vetores e Cotas
\draw[->, >=stealth, line width=1.5pt, color=black] (1,0.5) -- (1,2) node[above, fill=white, inner sep=2pt] {\textbf{F (Bíceps)}};
\draw[->, >=stealth, line width=1.5pt, color=black] (5,0.2) -- (5,-1.5) node[below, fill=white, inner sep=2pt] {\textbf{40 N}};

\draw[<->, >=stealth, dashed] (0, -1) -- (1, -1) node[midway, fill=white, inner sep=1pt] {\textbf{4 cm}};
\draw[<->, >=stealth, dashed] (1, -1) -- (5, -1) node[midway, fill=white, inner sep=1pt] {\textbf{36 cm}};
\draw[dashed, color=gray!80] (0,-0.3) -- (0,-1.2);
\draw[dashed, color=gray!80] (1,-0.5) -- (1,-1.2);
\draw[dashed, color=gray!80] (5,-0.5) -- (5,-1.2);
\end{tikzpicture}
\end{adjustbox}
\end{center}
Neste arranjo, o antebraço atua como uma alavanca interpotente. Calcule a força \textbf{F} que o músculo bíceps deve exercer para equilibrar o peso de \textbf{40 N} sustentado pela mão.
\par\vspace{\espacoQuestao}

\noindent\textcolor{CorLinha}{\rule{\linewidth}{0.5pt}} 
\par\vspace{\espacoPreQuestao}
\subsubsection*{22.} A invenção de ferramentas especializadas visou contornar limitações físicas do ser humano. Analise as afirmativas abaixo sobre ferramentas comuns e assinale Verdadeiro (\textbf{V}) ou Falso (\textbf{F}):
\begin{enumerate}[label=\Roman*., itemsep=\espacoAlt]
    \item Uma pinça de depilação é uma alavanca interpotente, projetada não para multiplicar a força, mas sim para fornecer controle fino e precisão na movimentação.
    \item O quebra-nozes atua como uma alavanca inter-resistente, pois a noz (resistência) é esmagada entre o eixo (dobradiça) e a força aplicada pelas mãos.
    \item Um alicate de corte é ineficiente para partir fios de aço espessos, pois atua como uma alavanca interpotente que divide a força aplicada pelo trabalhador.
\end{enumerate}
\par\vspace{\espacoQuestao}

\noindent\textcolor{CorLinha}{\rule{\linewidth}{0.5pt}} 
\par\vspace{\espacoPreQuestao}
\subsubsection*{23.} Para garantir a acessibilidade em calçadas, a norma brasileira NBR 9050 exige que rampas para cadeirantes não excedam uma proporção física limite para que o esforço não seja excessivo. A figura representa o projeto de uma dessas rampas.
\begin{center}
\begin{adjustbox}{trim=0cm 0cm 0cm 0cm, clip, max width=\larguraTikz}
\begin{tikzpicture}
% --- APLICANDO DROP SHADOW, CORES HEX E BORDAS ARREDONDADAS ---
\definecolor{rampaconcreto}{HTML}{BDC3C7}
\definecolor{piso}{HTML}{7F8C8D}

% LAYER 1: Chao e rampa
\draw[thick, color=piso] (-1,0) -- (9,0);
\draw[fill=rampaconcreto, drop shadow={opacity=0.15}, draw=gray!70] (0,0) -- (8,0) -- (8,0.8) -- cycle;

% LAYER 2: Cadeirante simplificado (Icone)
\draw[fill=white, thick] (2, 0.45) circle (0.2); 
\draw[thick, rounded corners=1pt] (2,0.45) -- (2, 0.9) -- (2.3, 0.9); 
\fill (2.1, 1.1) circle (0.1); 

% LAYER 3: Cotas com máscara
\draw[<->, >=stealth, thick] (8.3, 0) -- (8.3, 0.8) node[midway, right, fill=white, inner sep=1pt] {\textbf{0,6 m (Altura)}};
\draw[<->, >=stealth, thick] (0, -0.5) -- (8, -0.5) node[midway, fill=white, inner sep=2pt] {\textbf{D (Comprimento horizontal)}};
\draw[dashed, color=gray!80] (0,0) -- (0,-0.7);
\draw[dashed, color=gray!80] (8,0) -- (8,-0.7);
\end{tikzpicture}
\end{adjustbox}
\end{center}
Se a legislação local estabelece que a vantagem mecânica ideal (desprezando o atrito) deve exigir que a força aplicada seja de apenas $ \mfrac{1}{10} $ do peso total da cadeira de rodas e seu ocupante, qual deve ser o comprimento \textbf{D} aproximado da rampa para vencer a altura de \textbf{0,6 m}?
\par\vspace{\espacoQuestao}

\noindent\textcolor{CorLinha}{\rule{\linewidth}{0.5pt}} 
\par\vspace{\espacoPreQuestao}
\subsubsection*{24.} (Estilo ENEM - Oficial Adaptada) Um ciclista experiente sabe que o sistema de transmissão de sua bicicleta (coroa, corrente e catraca) funciona como um conjunto de roldanas acopladas por uma correia. Ao encarar uma subida muito íngreme, o ciclista decide acionar o câmbio para tornar a pedalada mais "leve". Do ponto de vista das máquinas simples, a pedalada se torna mais leve quando a corrente passa a ligar:
\begin{enumerate}[label=\alph*), itemsep=\espacoAlt]
    \item Uma coroa grande dianteira a uma catraca pequena traseira, garantindo maior velocidade da roda.
    \item Uma coroa pequena dianteira a uma catraca grande traseira, maximizando a vantagem mecânica de torque.
    \item Polias de raios exatamente iguais, cancelando o atrito da corrente.
    \item O eixo central diretamente à roda traseira, eliminando o uso da corrente transmissora.
    \item Uma catraca traseira cujo raio seja menor que o eixo central da roda, invertendo a força.
\end{enumerate}
\par\vspace{\espacoQuestao}

\noindent\textcolor{CorLinha}{\rule{\linewidth}{0.5pt}} 
\par\vspace{\espacoPreQuestao}
\subsubsection*{25.} Chegamos ao último desafio do seu treinamento em mecânica básica. Um gerador de energia pesando \textbf{2000 N} quebrou e precisa ser colocado em cima da carroceria de uma picape. No entanto, há apenas um funcionário disponível (que consegue exercer uma força máxima de apenas \textbf{250 N}), algumas tábuas de madeira lisas e cordas com polias à disposição.
\begin{boxresponda}
\textbf{Responda:} Proponha e descreva uma solução prática combinando DUAS máquinas simples diferentes que permitam a essa única pessoa erguer o gerador para a picape. Prove matematicamente que o esforço final será menor ou igual a 250 N.
\end{boxresponda}
\par\vspace{\espacoQuestao}

\end{multicols*}
\end{document}